% Part: turing-machines
% Chapter: undecidability
% Section: trakhtenbrot

\documentclass[../../../include/open-logic-section]{subfiles}

\begin{document}

\olfileid{tur}{und}{tra}
\olsection{Trakhtenbrot 定理}

\begin{explain}
  在 \olref[rep]{sec} 中，我们为图灵机~$M$ 和输入字符串~$w$ 定义了语句 $!T(M,w)$ 和~$!E(M,w)$。
  随后我们在 \olref[ver]{lem:valid-if-halt} 和 \olref[ver]{lem:halt-if-valid} 中证明了，$!T(M,w) \lif !E(M,w)$ 有效当且仅当 $M$ 以输入~$w$ 启动后最终停机。
  由于停机问题是不可判定的，这意味着一阶逻辑语句的有效性与可满足性是不可判定的（\cref{tur:und:uns:thm:decision-prob,tur:und:uns:cor:undecidable-sat}）。

  但是，语句的有效性和可满足性是针对任意结构（有限或无穷）定义的。
  你可能会想，判定一个语句在有限结构中是否可满足（或在所有有限结构中有效）或许会更容易。
  我们可以修改判定问题不可解性的证明，以表明情况并非如此。

  首先，如果你回顾 \olref[ver]{lem:halt-if-valid} 的证明，
  你会发现我们在那里构造了一个~$!T(M,w)$ 的模型~$\Struct M$，它精确地描述了机器~$M$ 以输入~$w$ 启动时的行为。
  该模型的论域是~$\Nat$，因而是无穷的。
  但如果 $M$ 确实以输入~$w$ 启动后停机，我们就可以用同样的方式构造一个有限模型~$\Struct M'$。
  假设 $M$ 以输入~$w$ 启动后在~$k$ 步停机。
  取论域~$\Domain{M'}$ 为集合 $\{0, \dots, n\}$，其中 $n$ 是~$k$ 和~$w$ 的长度中的较大者，并令
  \[
    \Assign{\prime}{M'}(x) = 
    \begin{cases}
      x + 1 &\text{若 $x < n$}\\
      n &\text{否则，}
    \end{cases}
  \]
  且令 $\tuple{x,y} \in \Assign{<}{M'}$ 成立当且仅当 $x < y$ 或 $x = y = n$。
  除此之外，$\Struct{M'}$ 的定义与~$\Struct{M}$ 完全相同。
  根据~$\Struct{M'}$ 的定义，正如在 \olref[ver]{lem:halt-if-valid} 的证明中那样，$\Sat{M'}{!T(M,w)}$。
  并且由于我们假设 $M$ 以输入~$w$ 启动后停机，$\Sat{M'}{!E(M,w)}$。
  因此，$\Struct{M'}$ 是~$!T(M,w) \land !E(M,w)$ 的一个有限模型（注意我们用~$\land$ 替代了 $\lif$）。
  
  我们的证明已经完成了一半：我们证明了如果 $M$ 以输入~$w$ 启动后停机，那么 $!T(M,w) \land !E(M,w)$ 有一个有限模型。
  不幸的是，其逆命题并不成立，即存在某些图灵机 $M$ 和输入~$w$，使得 $M$ 在~$w$ 上不停机，但 $!T(M,w)
  \land !E(M,w)$ 仍然有有限模型。
  例如，考虑只有状态 $q_0$ 且指令为 $\delta(q_0,\TMblank) = \tuple{q_0,\TMblank,\TMstay}$ 的机器~$M$。
  在空输入~$w = \emptyseq$ 上启动时，这台机器永远不会停机：它陷入了无限循环，但既不改变纸带也不移动读写头。
  所有的格局都是相同的（相同状态、相同读写头位置、相同纸带内容）。
  我们可以定义一个满足 $!T(M,\emptyseq) \land !E(M,\emptyseq)$ 的有限结构~$\Struct{M''}$（练习）。
  然而，我们可以通过适当地修改~$!T(M,w)$，以排除这样的结构。

  \begin{prob}
    设 $M$ 是一台只有状态 $q_0$ 和单条指令 $\delta(q_0,\TMblank) = \tuple{q,\TMblank,\TMstay}$ 的图灵机。
    令 $\Domain{M''} = \{0, 1, 2\}$、$\Assign{\prime}{M''}(0) =
    \Assign{\prime}{M'}(1) = 1$ 和 $\Assign{\prime}{M''}(2) = 2$，以及 $\Assign{<}{M''} = \{\tuple{0,1}, \tuple{1,1}, \tuple{2,2}\}$。
    定义 $\Assign{\Obj Q_{q_0}}{M''}$、$\Assign{\Obj
    S_{\TMblank}}{M''}$ 和 $\Assign{\Obj S_{\TMendtape}}{M''}$，使得 $!T(M,\emptyseq)$ 和 $!E(M,\emptyseq)$ 为真，并解释其原因。
    提示：注意 $\delta(q_0, \TMendtape)$ 是无定义的。
    确保
    \begin{align*}
    & \Obj Q_{q_0}(\num{1}, \num{n}) \land \Obj S_{\TMendtape}(\num{0}, \num{n})
      \land 
      \lforall[x][(\num{0} < x \lif \Obj S_\TMblank(x, \num{n}))] \text{\quad 对所有 $n \in \Nat$}\\
    & \lexists[y][(\Obj Q_{q_0}(\num{0}, y) \land \Obj S_{\TMendtape}(\num{0}, y))]
    \end{align*}
    在~$\Struct{M''}$ 中都为真。
  \end{prob}
\end{explain}

考虑描述图灵机~$M$ 在输入 $w = \sigma_{i_1}\dots\sigma_{i_k}$ 上操作的语句：
\begin{enumerate}
\item 描述数和~$<$ 的公理（与 \olref[rep]{sec} 中定义~$!T(M,w)$ 时相同）。
\item 描述初始格局的公理：与定义~$!T(M,w)$ 时相同。
\item 描述从一个格局转移到下一个格局的公理：

对于以下内容，令 $!A(x, y)$ 如前所述，并令
\[
  !B(y) \ident \lforall[x][(x < y \lif \eq/[x][y])].
\]
\begin{enumerate}
\item \ollabel{rep-right} 对于每条指令 $\delta(q_i, \sigma) =
  \tuple{q_j, \sigma', \TMright}$，语句：
\begin{align*}
& \lforall[x][\lforall[y][(
   (\Obj Q_{q_i}(x, y) \land \Obj S_{\sigma}(x, y)) \lif {}]] \\
&\qquad   (\Obj Q_{q_j}(x', y') \land \Obj S_{\sigma'}(x, y') \land
!A(x, y) \land !B(y')))
\end{align*}
\item \ollabel{rep-left} 对于每条指令 $\delta(q_i, \sigma) =
  \tuple{q_j, \sigma', \TMleft}$，语句：
\begin{align*}
& \lforall[x][\lforall[y][
    ((\Obj Q_{q_i}(x', y) \land \Obj S_{\sigma}(x', y)) \lif {}]]\\
& \qquad   (\Obj Q_{q_j}(x, y') \land \Obj S_{\sigma'}(x', y') \land
!A(x, y))) \land {}\\
& \lforall[y][((\Obj Q_{q_i}(\num{0}, y) \land \Obj S_{\sigma}(\num{0},
    y)) \lif {}]\\
& \qquad (\Obj Q_{q_j}(\num{0}, y') \land \Obj S_{\sigma'}(\num{0},
  y') \land !A(\num{0}, y) \land !B(y')))
\end{align*}
\item \ollabel{rep-stay} 对于每条指令 $\delta(q_i, \sigma) =
  \tuple{q_j, \sigma', \TMstay}$，语句：
\begin{align*}
& \lforall[x][\lforall[y][(
   (\Obj Q_{q_i}(x, y) \land \Obj S_{\sigma}(x, y)) \lif {}]] \\
&\qquad (\Obj Q_{q_j}(x, y') \land \Obj S_{\sigma'}(x, y') \land
!A(x, y) \land !B(y')))
\end{align*}
\end{enumerate}
如你所见，描述~$M$ 转移的语句与~$!T(M,w)$ 中对应的语句相同，只是我们在末尾添加了 $!B(y')$。
$!B(y')$~确保了“下一个”格局的编号 $y'$ 与之前的所有编号 $\Obj 0$, $\Obj 0'$, …… 都不同。
% \item Sentences that express consistency conditions:
% \begin{enumerate}
%   \item\ollabel{rep-con-symb}%
%   !!^a{sentence} that says that no square can have more than one
%   symbol on it at any given time:
%   \[\lforall[x][\lforall[y][(\Obj S_{\sigma}(x, y) \lif \lnot \Obj
%   S_{\sigma'}(x, y))]],\]
%   for every pair $\sigma \neq \sigma'$ of tape symbols of~$T$.
%   \item\ollabel{rep-con-state}%
%   !!^a{sentence} that says that $M$ cannot be in more than one state
%   at any given time:
%   \[\lforall[x][\lforall[y][(\Obj Q_{q}(x, y) \lif
%   \lforall[z][\lnot \Obj
%   Q_{q'}(z, y))]],\]
%   for every pair $q \neq q'$ of states of~$T$.
%   \item\ollabel{rep-con-tape}%
%   !!^a{sentence} that says that $M$ cannot be on more than one tape
%   square at any given time:
%   \[\lforall[x][\lforall[y][(\Obj Q_{q}(x, y) \lif
%   \lforall[z][\eq/[z][y]\lif \lnot \Obj
%   Q_{q'}(z, y))]],\]
%   for every pair $q$, $q'$ of states of~$T$.
% \end{enumerate}
\end{enumerate}
令 $!T'(M, w)$ 为上述所有关于图灵机~$M$ 和输入~$w$ 的语句的合取。

\begin{lem}\ollabel{lem:halts-sat}
  若 $M$ 以输入~$w$ 启动后停机，则 $!T'(M,w) \land !E(M,w)$ 有一个有限模型。
\end{lem}

\begin{proof}
  令 $\Struct{M'}$ 如 \olref[ver]{lem:halt-if-valid} 的证明中所述，除以下不同外：\begin{align*}
    \Domain{M'} & = \{0, \dots, n\},\\
    \Assign{\prime}{M'}(x) & = 
    \begin{cases}
      x + 1 &\text{若 $x < n$}\\
      n &\text{否则，}
    \end{cases} \\
    \tuple{x,y} \in \Assign{<}{M'} &\text{当且仅当 $x < y$ 或 $x = y = n$，}
  \end{align*}
  其中 $n = \max(k,\len{w})$，且 $k$~是使得 $M$ 以输入~$w$ 启动后在第~$k$ 步停机的最小数。我们将 $\Sat{M'}{!T'(M,w) \land E(M,w)}$ 的验证留作练习。
\end{proof}

\begin{prob}
  请完成 \olref[tur][und][tra]{lem:halts-sat} 的证明，即证明 $\Sat{M'}{!T(M,w) \land E(M,w)}$。
\end{prob}

\begin{lem}\ollabel{lem:sat-halts}
  如果 $!T'(M,w) \land !E(M,w)$ 有一个有限模型，那么 $M$ 以输入~$w$ 启动后会停机。
\end{lem}

\begin{proof}
  我们证明其逆否命题。假设 $M$ 以输入~$w$ 启动后不停机。如果 $!T'(M,w) \land !E(M,w)$ 根本没有模型，则得证。
  所以假设 $\Struct{M}$ 是~$!T(M,w) \land !E(M, w)$ 的一个模型。我们需要证明它不可能是有限的。
  
  正如在 \olref[ver]{lem:config} 中那样，我们可以证明：如果 $M$ 以输入~$w$ 启动后，在~$n$ 步之后仍未停机，那么 $!T'(M,w)
  \Entails !C(M, w, n) \land !B(\num{n})$。由于 $M$ 以输入~$w$ 启动后不停机，因此对所有 $n \in \Nat$，$!T'(M,w) \Entails !C(M, w, n) \land
  !B(\num{n})$。注意，根据 \olref[rep]{prop:mlessk}，对所有 $k < n$，$!T'(M,w) \Entails \num{k} < \num{n}$。另外 $!B(\num{n}) \Entails \num{k} < \num{n} \lif
  \eq/[\num{k}][\num{n}]$。所以，对所有 $k < n$，$\Sat{M}{\eq/[\num{k}][\num{n}]}$，即无穷多个项~$\num{k}$ 在~$\Struct{M}$ 中必须取不同的值。但这要求 $\Domain{M}$ 是无穷的，所以 $\Struct{M}$ 不可能是~$!T'(M,w) \land !E(M, w)$ 的有限模型。
\end{proof}

\begin{prob}
  请完成 \olref[tur][und][tra]{lem:sat-halts} 的证明，即证明：如果 $M$ 以输入~$w$ 启动后，在~$n$ 步之后仍未停机，那么 $!T'(M,w) \Entails !B(\num{n})$。
\end{prob}

\begin{thm}[Trakhtenbrot 定理]
  \ollabel{thm:trakhtenbrodt}
  任意一阶逻辑语句是否拥有有限模型（即是否有限可满足），这是不可判定的。
\end{thm}

\begin{proof}
  假设存在一台图灵机~$F$ 可以判定有限可满足性问题。那么给定任意图灵机~$M$ 和输入~$w$，我们都可以计算出语句~$!T'(M,w) \land !E(M,w)$，并用~$F$ 来判定它是否有一个有限模型。根据 \cref{tur:und:tra:lem:halts-sat,tur:und:tra:lem:sat-halts}，它有有限模型当且仅当 $M$ 以输入~$w$ 启动后停机。这样我们就可以用 $F$ 来解决停机问题，但我们知道停机问题是不可解的。
\end{proof}

\begin{cor}\ollabel{cor:fproof-incomp}%
  不可能存在一个关于\emph{有限}有效性的既可靠又完备的推导系统，即，不存在这样一个推导系统，使得对于每个有限结构~$\Struct{M}$，$\Proves !B$ 当且仅当 $\Sat{M}{!B}$。
\end{cor}

\begin{proof}
  留作练习。
\end{proof}

\begin{prob}
  证明 \olref[tur][und][tra]{cor:fproof-incomp}。注意，$!B$ 在每个有限结构中可满足当且仅当 $\lnot !B$ 不是有限可满足的。请解释为何有限可满足性在 \olref[tur][und][uns]{thm:valid-ce} 的意义上是半可判定的。并由此论证：如果存在一个关于有限有效性的推导系统，那么有限可满足性就是可判定的。
\end{prob}


\end{document}